\documentclass[12pt]{article}
\usepackage[utf8]{inputenc}
\usepackage[T1]{fontenc}
\usepackage[margin=1in]{geometry}
\usepackage{setspace}
\usepackage{xcolor}
\usepackage[colorlinks=true, urlcolor=blue!60!black, linkcolor=blue!60!black]{hyperref}
\setlength{\parindent}{0em}
\setlength{\parskip}{0.45em}
\renewcommand{\baselinestretch}{1.0}
\begin{document}
\begin{center}
{\Large \textit{Cooper v. Aaron}, 358 U.S. 1 (1958)\par}
\vspace{0.3em}
{\normalsize Case Brief --- One-Page Edition\par}
\end{center}

\textbf{Memory Jogger.} Little Rock: official defiance cannot suspend desegregation, and the Court's reading of the Constitution binds the States (\href{https://supreme.justia.com/cases/federal/us/358/1/}{Justia}).

\textbf{Facts.} Little Rock adopted a gradual desegregation plan. Governor Orval Faubus used the National Guard to block the ``Little Rock Nine,'' a mob forced their removal, and President Eisenhower sent Army troops.

\textbf{Procedural History.} Citing ``chaos, bedlam and turmoil,'' the School Board won a suspension of the plan; the Eighth Circuit reversed, and the Supreme Court affirmed in a special term.

\textbf{Votes.} 9--0, in an opinion signed by all nine Justices; Frankfurter concurred.

\textbf{Holding.} School officials act for the State and cannot suspend rights because of disorder provoked by the Governor and Legislature. \textit{Brown}'s interpretation is the supreme law of the land, binding all state officials.

\textbf{Reasoning.} The turmoil was ``directly traceable'' to state officials; ``law and order are not here to be preserved by depriving the Negro children of their constitutional rights.'' Reading \textit{Marbury} with Article VI, officials sworn to support the Constitution cannot ``war against'' it. Frankfurter warned that yielding would mean ``that law should bow to force.''

\textbf{Examples.} \emph{Same-Side:} a county clerk refusing marriage licenses. \emph{Opposite-Side:} a presidential veto of a bill the Court would uphold. \emph{Fence-Sitter:} a statute enacted to test a precedent.

\textbf{Critique.} Article VI makes the Constitution, not the Court's opinions, supreme; the Lincoln and \textit{Dred Scott} objection.

\textbf{Key Quotation.} ``This decision declared the basic principle that the federal judiciary is supreme in the exposition of the law of the Constitution, and that principle has ever since been respected by this Court and the Country as a permanent and indispensable feature of our constitutional system.'' --- Opinion of the Court, 358 U.S. at 18.

\end{document}
